\section{Rotations and limitations}\label{sec:examples}
The following exact examples delineate the scope of the preceding limit
theorems. Throughout this section, the planar rotation is
\[
 U_\theta=\begin{pmatrix}
 \cos\theta&-\sin\theta\\
 \sin\theta&\cos\theta
 \end{pmatrix}.
\]

\begin{proposition}[An exactly solvable rotating reaction family]\label{prop:reactionrotation}
Use entrywise divergence on a unit-volume domain, uniform endpoints
$G_0=I_2\,dx$, $G_1=\diag(2,1)\,dx$, and any fixed PSD transport weight. Put
$H_{\varepsilon,\theta}=U_\theta\diag(1,\varepsilon)U_\theta^T$,
where $U_\theta$ is the planar rotation, $\varepsilon>0$. The singular target
$H_* =\diag(1,0)$ has value $V_*=2(\sqrt2-1)^2$. Writing
$p=\sin^2\theta$ and $Z=1+p+2\sqrt2\varepsilon^2+(2-p)\varepsilon^4$,
the current value is
\begin{equation}\label{eq:rotationformula}
 V=2(3-p)+\frac{2(p/\varepsilon)^2}{(\sqrt{1+p}+1)^2}
 -\frac{4[2\sqrt2+(2-p)\varepsilon^2]}{\sqrt Z+\sqrt{1+p}}.
\end{equation}
As $\varepsilon\to0$ and $\theta\to0$, an extended limit
$\sin^2\theta/\varepsilon\to\ell$ implies $V\to V_*+\ell^2/2$.
In particular $\theta=o(\sqrt\varepsilon)$ gives $V\to V_*$,
$\theta/\sqrt\varepsilon\to a$ gives $V\to V_*+a^4/2$, and
$|\theta|/\sqrt\varepsilon\to\infty$ gives
$V\sim\theta^4/(2\varepsilon^2)\to\infty$.
\end{proposition}
\begin{proof}
The factor construction in Lemma~\ref{lem:pure}, applied to the uniform
endpoints, gives an admissible zero-flux path with action
$2d_B^2(X_0,X_1)$, where $A_0=I_2$, $A_1=\diag(2,1)$ and
$X_i=H^{-1}A_iH^{-1}$. Its space-constant
Riccati certificates give the matching lower bound against all competitors.
Thus this is the original minimum for any transport weight.

If the eigenvalues of a two-dimensional PSD matrix $Y$ are
$\lambda_1,\lambda_2\ge0$, then
$(\tr\sqrt Y)^2=(\sqrt{\lambda_1}+\sqrt{\lambda_2})^2
=\tr Y+2\sqrt{\det Y}$. In the eigenbasis of $H$, one has
\[
 U_\theta^TA_1U_\theta=
 \begin{pmatrix}2-p&-\sin\theta\cos\theta\\
 -\sin\theta\cos\theta&1+p\end{pmatrix},\qquad
 U_\theta^TX_0U_\theta=\diag(1,\varepsilon^{-2}).
\]
Congruence of the first matrix by $\diag(1,\varepsilon^{-1})$ gives
$U_\theta^TX_1U_\theta$. Taking traces and determinants yields
\begin{align*}
 \tr X_0+\tr X_1&=3-p+(2+p)/\varepsilon^2,\\
 \tr(X_0^{1/2}X_1X_0^{1/2})&=2-p+(1+p)/\varepsilon^4,\\
 \det(X_0^{1/2}X_1X_0^{1/2})&=2/\varepsilon^4.
\end{align*}
The square-root trace is $\sqrt Z/\varepsilon^2$, so substitution gives
\[
 V=2(3-p)+\frac{2}{\varepsilon^2}(2+p-2\sqrt Z).
\]
To expose the cancellation, first split $\sqrt Z$ at $\sqrt{1+p}$ and use
\[
 \sqrt Z-\sqrt{1+p}
 =\frac{2\sqrt2\varepsilon^2+(2-p)\varepsilon^4}
        {\sqrt Z+\sqrt{1+p}},\qquad
 2+p-2\sqrt{1+p}
 =\frac{p^2}{(\sqrt{1+p}+1)^2}.
\]
These two identities give \eqref{eq:rotationformula}. Its first and last
terms together tend to $6-4\sqrt2=V_*$ and remain bounded, since $p\to0$
and both square roots in the last denominator tend to one. The middle
term tends to $\ell^2/2$ if $p/\varepsilon\to\ell<\infty$ and is
asymptotic to $p^2/(2\varepsilon^2)$ if that ratio diverges. In the latter
case it dominates the bounded terms. Finally
$\sin\theta/\theta\to1$ converts these conclusions into each of the stated
thresholds in $\theta$.
For the singular target the uniform density
$M(t)=\diag((1+(\sqrt2-1)t)^2,1)$, with
$G=M(t)\,dt\,dx$, $R=\dot M(t)\,dt\,dx$ and $q=0$, is admissible by
integration by parts in time. Put $a=\sqrt2-1$. Its sole nonzero reaction
entry is $2a(1+at)$ in the supported first column. To compute the cost,
maximize first over $A$ by setting
$A=-BLL^TB^T/2-CH_*^2C^T/2$. The zero flux makes $B=0$ optimal.
Writing $m=(1+at)^2$ and $z=2a(1+at)$, the remaining objective is
\[
 zC_{11}-\tfrac12m C_{11}^2-\tfrac12C_{21}^2
 \le\frac{z^2}{2m}=2a^2.
\]
Equality holds at $C_{11}=z/m$, $C_{21}=0$, with the unused second column
set to zero. The unit-volume domain therefore gives total action $V_*$. For the lower bound, $I_2\succeq H_*^2$ and
Lemma~\ref{lem:comparison} imply that the identity reaction weight is
no more costly than $H_*$. The latter minimum is thus at least the
identity-weight Bures value $V_*$. This proves equality at the singular target.
\end{proof}

\begin{proposition}[A rotating transport jump]\label{prop:transportrotation}
In the scalar unit-disk problem with ordinary divergence and $H=1$, take
$G_0=\delta_{(-1/2,0)}$, $G_1=\delta_{(1/2,0)}$ and
$L_\varepsilon=U_{\sqrt\varepsilon}\diag(1,\varepsilon)U_{\sqrt\varepsilon}^T$.
For sufficiently small $\varepsilon>0$, $V(L_\varepsilon,1)=4$, whereas
$V(\diag(1,0),1)\le1/2$.
\end{proposition}
\begin{proof}
Write $x_0=(-1/2,0)$ and $x_1=(1/2,0)$. A dual cutoff in the expensive
direction gives the lower bound; two stationary reaction paths attain it.
Let $u=(\sin\sqrt\varepsilon,-\cos\sqrt\varepsilon)$, the expensive unit
eigendirection. Define the globally $C^1$ function
\[
 h(z)=\begin{cases}0&z\le-\pi/2,\\(1+\sin z)/2&|z|<\pi/2,\\1&z\ge\pi/2,
 \end{cases}\qquad f(x)=h(u\cdot x/\varepsilon).
\]
At both breakpoints the interior derivative $\cos z/2$ vanishes, matching
the constant pieces; hence $h$ is $C^1$. Moreover $(h')^2=h(1-h)$ on each
piece and at the breakpoints. Since $L_\varepsilon u=\varepsilon u$,
\[
 |L_\varepsilon\nabla f|^2
 =|L_\varepsilon u|^2\varepsilon^{-2}(h'(u\cdot x/\varepsilon))^2
 =f(1-f).
\]
The values $u\cdot x_0=-\sin\sqrt\varepsilon/2$ and
$u\cdot x_1=\sin\sqrt\varepsilon/2$ lie on opposite constant plateaus
once $\sin\sqrt\varepsilon\ge\pi\varepsilon$. This holds eventually,
because $\sin\sqrt\varepsilon/\varepsilon\to\infty$, and gives
$f(x_0)=0$, $f(x_1)=1$.
For $\delta>0$ set $a(t)=-2/(1+\delta-t)$, $b(t)=2/(t+\delta)$,
$\Phi=(1-f)a+fb$. Direct differentiation gives $a'=-a^2/2$ and
$b'=-b^2/2$. Also
\[
 \Phi^2=(1-f)a^2+fb^2-f(1-f)(b-a)^2,\qquad
 |L_\varepsilon\nabla\Phi|^2=(b-a)^2f(1-f).
\]
Adding half of these two identities to
$\partial_t\Phi=-((1-f)a^2+fb^2)/2$ proves
$\partial_t\Phi+|L_\varepsilon\nabla\Phi|^2/2+\Phi^2/2=0$.
Its endpoint pairing is $\Phi(1,x_1)-\Phi(0,x_0)=b(1)-a(0)=4/(1+\delta)$. Fenchel inequality and
\eqref{eq:ce}, followed by $\delta\downarrow0$, give the lower bound $4$.
The admissible zero-flux path
$G_t=(1-t)^2\delta_{x_0}+t^2\delta_{x_1}$,
$R_t=-2(1-t)\delta_{x_0}+2t\delta_{x_1}$ has action $4$.
Indeed, its reaction is the time derivative of its mass, which verifies
\eqref{eq:ce}. Using the reference measure
$dt\otimes(\delta_{x_0}+\delta_{x_1})$, each atom with scalar mass $m>0$
and reaction $z$ has support cost
\[
 \sup_c\{zc-mc^2/2\}=z^2/(2m),
\]
because the zero flux makes the transport dual zero optimal, the mass
constraint is saturated, and the maximum occurs at $c=z/m$.
Thus the cost is $[-2(1-t)]^2/[2(1-t)^2]=2$ at the first atom for $t<1$
and $(2t)^2/(2t^2)=2$ at the second for $t>0$. At the two exceptional endpoint
times the measures have zero mass, so integrating gives $2+2=4$.
At the singular target, the horizontal moving unit atom of velocity $(1,0)$
has zero reaction and action at most $1/2$. Explicitly, push $ds$ forward
by $s\mapsto(s,(-1/2+s,0))$ to form $G$ and take $q=(1,0)G$, $R=0$.
The chain rule and interval fundamental theorem give \eqref{eq:ce} for every
closed-domain test. For a scalar admissible dual triple,
$A+B_1\le-B_1^2/2+B_1\le1/2$, which proves the action bound.
\end{proof}

\begin{proposition}[Frobenius divergence does not ensure a value limit]\label{prop:norm}
For scalar ordinary divergence on the unit disk, $H=1$, and endpoints
$\delta_{(0,-1)}$, $\delta_{(0,1)}$, an alternating sequence
$L_{2j}=\diag(t_{2j},1)$, $L_{2j+1}=\diag(1,t_{2j+1})$, $0<t_j\to\infty$,
has $\norm{L_j}_F^2\to\infty$ but nonconvergent minimum values.
\end{proposition}
\begin{proof}
For the even sequence $\Phi(t,x,y)=y/4-t/8$ satisfies
$|\Phi|\le3/8$ on the time--disk cylinder and
$L_{2j}\nabla\Phi=(0,1/4)$, so it has
\[
 \partial_t\Phi+\tfrac12|L_{2j}\nabla\Phi|^2+\tfrac12\Phi^2
 \le-\tfrac18+\tfrac1{32}+\tfrac9{128}=-\tfrac3{128}.
\]
Its endpoint pairing is
$\Phi(1,0,1)-\Phi(0,0,-1)=1/8+1/4=3/8$, so the even minima are at least
$3/8$ by \eqref{eq:ce} and the support inequality.
For the odd sequence use $\gamma(s)=(0,2s-1)$,
$G=(s\mapsto(s,\gamma(s)))_\#ds$, $q=(0,2)G$, $R=0$.
The trajectory lies in the interior for interior times, so the chain rule
and interval fundamental theorem verify every original weak test.
The support constraint gives
$A+2B_2\le-t_{2j+1}^2B_2^2/2+2B_2\le2/t_{2j+1}^2$ eventually.
Hence the balanced action is at most $2/t_{2j+1}^2\to0$.
Thus the values cannot converge.
\end{proof}

\begin{proposition}[The operator survives uniformly cheap transport]\label{prop:operator}
On the unit disk, for the Kantorovich--Bures operator
$D^*\Phi=\operatorname{div}\Phi$,
$n=d=2$, $k=1$, $H=I_2$, and endpoints
$I_2\delta_{(0,-1)}$, $I_2\delta_{(0,1)}$, one has $V(tI_1,I_2)\ge1/4$
for all $t>0$, although the total endpoint matrices coincide.
\end{proposition}
\begin{proof}
Use $\Phi(s,x,y)=\diag(y/4-s/8,-s/8)$. Its row divergence is
$(\partial_x\Phi_{11}+\partial_y\Phi_{12},
  \partial_x\Phi_{21}+\partial_y\Phi_{22})=(0,0)$.
Here $s\in[0,1]$ is path time. On the unit disk $|y|\le1$, hence
$|y/4-s/8|\le3/8$ and $|s/8|\le1/8$. The diagonal entries of
$\partial_s\Phi+\Phi^2/2$ are therefore at most
$-1/8+9/128\le-3/128$ and $-1/8+1/128=-15/128$.
At the terminal point $(0,1)$ the trace of $\Phi(1)$ is zero; at the
initial point $(0,-1)$ the trace of $\Phi(0)$ is $-1/4$. The endpoint
pairing is therefore $1/4$. Thus the dual lower bound is independent
of the scalar transport weight. The Bures expression of the total matrices
is zero, so the entrywise-divergence identification cannot extend to this
operator.
\end{proof}

\begin{remark}\label{rem:nocompact}
Unscaled flux need not be precompact even for vanishing-gap paths in the
large-transport regime. On a unit spatial interval, let $s\in[0,1]$ denote
path time and $\tau>0$ the transport scale. Take equal endpoints $I_2\,dx$,
$H=I_2$, and $J=\left(\begin{smallmatrix}0&1\\-1&0\end{smallmatrix}\right)$.
For entrywise divergence the space--time measures
\[
 G^{(\tau)}=I_2\,ds\,dx,\qquad
 q^{(\tau)}=\sqrt\tau J\,ds\,dx,\qquad R^{(\tau)}=0
\]
have densities constant in path time. The skew matrix $J$ pairs to zero
with every symmetric spatial derivative, so these paths satisfy all
symmetric weak tests. Their action at weight $\tau I_2$ is
$\norm J_F^2/(2\tau)$, while flux variation is $\sqrt\tau\norm J_F$.
For this exact calculation, saturate the mass-dual constraint at $g=I_2$.
With zero reaction, $C=0$ is optimal, so the support cost is
\[
 \sup_B\{B\cdot z-\tfrac12\tau^2\norm B_F^2\}
 =\frac{\norm z_F^2}{2\tau^2},
\]
by completing the square, with equality at $B=z/\tau^2$. Inserting
$z=\sqrt\tau J$ and integrating over the unit time--space rectangle gives
the asserted action exactly.
The stationary zero-flux path has action zero, so the gaps vanish as
$\tau\to\infty$. Pairing the flux with the constant test selecting its
$(1,2)$ entry gives $\sqrt\tau$; thus no weak-star convergent subsequence
exists along $\tau\to\infty$. This explains why
Theorem~\ref{thm:uniform} asserts energy convergence without flux compactness.
\end{remark}
